% 与 ctexbook + XeLaTeX 配合；中文字体由本文件设置
\usepackage{geometry}
\geometry{a4paper, margin=2.5cm, headheight=14pt}

% Windows 默认；若缺字体请改成本机已安装的中文字体名
\setCJKmainfont{SimSun}
\setCJKsansfont{Microsoft YaHei}
\setCJKmonofont{FangSong}
\setmainfont{Times New Roman}
\setmonofont{Consolas}[Scale=0.92]

\usepackage{graphicx}
\usepackage{booktabs}
\usepackage{longtable}
\usepackage{tabularx}
\usepackage{multirow}
\usepackage{array}
\usepackage{enumitem}
\usepackage{xcolor}
\usepackage{fancyhdr}
\usepackage{titlesec}
\usepackage{caption}
\usepackage{listings}
\usepackage{tcolorbox}
\usepackage{amsmath}
\usepackage{microtype}
\usepackage{underscore}
\usepackage{hyperref}

\hypersetup{
  colorlinks=true,
  linkcolor=blue!50!black,
  urlcolor=blue!60!black,
  citecolor=green!40!black,
  bookmarksnumbered=true,
  pdftitle={PrimeTime User Guide 中文学习译本},
  pdfauthor={Unofficial Study Translation}
}

\pagestyle{fancy}
\fancyhf{}
\fancyhead[L]{\small PrimeTime UG T-2022.03}
\fancyhead[R]{\small 中文学习译本}
\fancyfoot[C]{\thepage}
\renewcommand{\headrulewidth}{0.4pt}

\titleformat{\chapter}[display]
  {\Huge\bfseries}{\chaptertitlename\ \thechapter}{12pt}{\Huge}
\titleformat{\section}{\Large\bfseries}{\thesection}{1em}{}
\titleformat{\subsection}{\large\bfseries}{\thesubsection}{1em}{}
\titleformat{\subsubsection}{\normalsize\bfseries}{\thesubsubsection}{1em}{}

\newcommand{\bichapter}[2]{%
  \chapter[#1]{#1\\[0.35em]{\LARGE\normalfont\itshape #2}}%
}

\tcbuselibrary{skins,breakable}
\newtcolorbox{noteBox}{
  colback=yellow!8, colframe=orange!60!black,
  fonttitle=\bfseries, title=注意,
  breakable, left=4pt, right=4pt
}
\newtcolorbox{seeAlsoBox}{
  colback=blue!4, colframe=blue!50!black,
  fonttitle=\bfseries, title=参见,
  breakable, left=4pt, right=4pt
}

\lstdefinestyle{tcl}{
  language=Tcl,
  basicstyle=\ttfamily\small,
  keywordstyle=\color{blue!70!black}\bfseries,
  commentstyle=\color{green!40!black}\itshape,
  stringstyle=\color{red!60!black},
  showstringspaces=false,
  breaklines=true,
  frame=single,
  rulecolor=\color{black!30},
  backgroundcolor=\color{gray!5},
  columns=fullflexible,
  keepspaces=true,
  tabsize=2,
  captionpos=b
}
\lstset{style=tcl}
\newcommand{\cmd}[1]{\texttt{#1}}
\newcommand{\opt}[1]{\texttt{#1}}
\newcommand{\appopt}[1]{\texttt{#1}}

\captionsetup{font=small, labelfont=bf}
\renewcommand{\figurename}{图}
\renewcommand{\tablename}{表}
\renewcommand{\lstlistingname}{代码}

\setlist[itemize]{leftmargin=1.8em, itemsep=0.25em, topsep=0.4em}
\setlist[enumerate]{leftmargin=1.8em, itemsep=0.25em, topsep=0.4em}

\newcommand{\copyrightnotice}{%
\begin{center}
\small
原文版权归 Synopsys, Inc.\ 所有（Version T-2022.03, March 2022）。\\
本文件为\textbf{非官方个人学习译本}，仅供学习参考，不得用于商业传播。\\
命令名、选项名、文件名等技术标识符保持英文原文不译。
\end{center}
}

\newcommand{\chapterstub}[2]{%
  \begin{tcolorbox}[colback=gray!5, colframe=gray!60, title=翻译占位]
  本章对应原书第~#1~章：\textit{#2}。\\
  正文待续译。可将原版 PDF 对应页提取后按第~1~章体例写入本文件。
  \end{tcolorbox}
}

\newcommand{\figplaceholder}[3]{%
\begin{figure}[htbp]
\centering
\fbox{\parbox{0.85\textwidth}{\centering\vspace{1.2cm}
\textit{（原书 #1）}\\[0.3em]
\small 请参阅原版 PDF 插图；可将截图放入 \texttt{figures/} 目录后用 \texttt{\textbackslash includegraphics} 替换。
\vspace{1.2cm}}}
\caption{#2}
\label{#3}
\end{figure}
}
